% =====================================================================
% Figura: carga puntual q frente a una esfera conductora a tierra
% (a = 1, d = 2.2). (a) Problema real: las líneas que llegan a la esfera
% marcan la carga inducida (fracción a/d del flujo de q); el resto escapa.
% (b) Problema equivalente: imagen q' = -q a/d en d' = a^2/d; dentro de la
% esfera las líneas son virtuales; la esfera es la equipotencial phi = 0.
% =====================================================================
\begin{tikzpicture}
  \def\esc{0.98}
  \def\caja{(-2.1*\esc,-2.35*\esc) rectangle (4.3*\esc,2.35*\esc)}
  % ---------------- (a)
  \begin{scope}
    \fill[panel] \caja;
    \begin{scope}
      \clip \caja;
      % Generado automáticamente por generar_figuras.py; no editar.
\draw[flinea, flecha=0.608cm] plot coordinates {(2.094,0.029) (2.058,0.045) (2.022,0.061) (1.987,0.077) (1.951,0.094) (1.915,0.110) (1.879,0.125) (1.843,0.141) (1.807,0.157) (1.771,0.172) (1.735,0.187) (1.699,0.202) (1.662,0.216) (1.626,0.230) (1.589,0.243) (1.552,0.256) (1.514,0.268) (1.477,0.279) (1.439,0.290) (1.401,0.299) (1.363,0.307) (1.324,0.314) (1.285,0.320) (1.246,0.324) (1.207,0.326) (1.168,0.327) (1.129,0.326) (1.090,0.323) (1.051,0.318) (1.012,0.311) (0.974,0.302) (0.936,0.291) (0.936,0.291)};
\draw[flinea, flecha=0.608cm] plot coordinates {(2.094,-0.029) (2.058,-0.045) (2.022,-0.061) (1.987,-0.077) (1.951,-0.094) (1.915,-0.110) (1.879,-0.125) (1.843,-0.141) (1.807,-0.157) (1.771,-0.172) (1.735,-0.187) (1.699,-0.202) (1.662,-0.216) (1.626,-0.230) (1.589,-0.243) (1.552,-0.256) (1.514,-0.268) (1.477,-0.279) (1.439,-0.290) (1.401,-0.299) (1.363,-0.307) (1.324,-0.314) (1.285,-0.320) (1.246,-0.324) (1.207,-0.326) (1.168,-0.327) (1.129,-0.326) (1.090,-0.323) (1.051,-0.318) (1.012,-0.311) (0.974,-0.302) (0.936,-0.291) (0.936,-0.291)};
\draw[flinea, flecha=0.751cm] plot coordinates {(2.106,0.047) (2.078,0.074) (2.049,0.101) (2.020,0.128) (1.992,0.154) (1.963,0.181) (1.934,0.207) (1.905,0.233) (1.875,0.259) (1.846,0.285) (1.816,0.311) (1.786,0.336) (1.755,0.360) (1.725,0.385) (1.693,0.408) (1.662,0.432) (1.630,0.454) (1.597,0.476) (1.564,0.497) (1.530,0.517) (1.496,0.535) (1.461,0.553) (1.425,0.570) (1.389,0.585) (1.352,0.598) (1.315,0.610) (1.277,0.620) (1.239,0.628) (1.200,0.634) (1.161,0.637) (1.122,0.638) (1.082,0.637) (1.043,0.632) (1.005,0.626) (0.967,0.616) (0.930,0.604) (0.893,0.589) (0.858,0.571) (0.825,0.551) (0.815,0.545)};
\draw[flinea, flecha=0.751cm] plot coordinates {(2.106,-0.047) (2.078,-0.074) (2.049,-0.101) (2.020,-0.128) (1.992,-0.154) (1.963,-0.181) (1.934,-0.207) (1.905,-0.233) (1.875,-0.259) (1.846,-0.285) (1.816,-0.311) (1.786,-0.336) (1.755,-0.360) (1.725,-0.385) (1.693,-0.408) (1.662,-0.432) (1.630,-0.454) (1.597,-0.476) (1.564,-0.497) (1.530,-0.517) (1.496,-0.535) (1.461,-0.553) (1.425,-0.570) (1.389,-0.585) (1.352,-0.598) (1.315,-0.610) (1.277,-0.620) (1.239,-0.628) (1.200,-0.634) (1.161,-0.637) (1.122,-0.638) (1.082,-0.637) (1.043,-0.632) (1.005,-0.626) (0.967,-0.616) (0.930,-0.604) (0.893,-0.589) (0.858,-0.571) (0.825,-0.551) (0.815,-0.545)};
\draw[flinea, flecha=0.971cm] plot coordinates {(2.119,0.057) (2.097,0.090) (2.076,0.123) (2.054,0.156) (2.032,0.188) (2.011,0.221) (1.989,0.253) (1.966,0.286) (1.944,0.318) (1.921,0.350) (1.898,0.382) (1.875,0.413) (1.851,0.444) (1.827,0.475) (1.802,0.505) (1.777,0.535) (1.751,0.565) (1.725,0.594) (1.698,0.623) (1.670,0.650) (1.642,0.678) (1.613,0.704) (1.584,0.730) (1.553,0.755) (1.522,0.779) (1.490,0.801) (1.458,0.823) (1.424,0.843) (1.390,0.863) (1.355,0.880) (1.319,0.896) (1.283,0.911) (1.246,0.923) (1.208,0.934) (1.170,0.943) (1.131,0.949) (1.092,0.954) (1.053,0.955) (1.014,0.955) (0.975,0.952) (0.936,0.946) (0.898,0.937) (0.860,0.925) (0.824,0.911) (0.788,0.894) (0.754,0.875) (0.722,0.853) (0.691,0.828) (0.662,0.802) (0.635,0.773) (0.622,0.757)};
\draw[flinea, flecha=0.971cm] plot coordinates {(2.119,-0.057) (2.097,-0.090) (2.076,-0.123) (2.054,-0.156) (2.032,-0.188) (2.011,-0.221) (1.989,-0.253) (1.966,-0.286) (1.944,-0.318) (1.921,-0.350) (1.898,-0.382) (1.875,-0.413) (1.851,-0.444) (1.827,-0.475) (1.802,-0.505) (1.777,-0.535) (1.751,-0.565) (1.725,-0.594) (1.698,-0.623) (1.670,-0.650) (1.642,-0.678) (1.613,-0.704) (1.584,-0.730) (1.553,-0.755) (1.522,-0.779) (1.490,-0.801) (1.458,-0.823) (1.424,-0.843) (1.390,-0.863) (1.355,-0.880) (1.319,-0.896) (1.283,-0.911) (1.246,-0.923) (1.208,-0.934) (1.170,-0.943) (1.131,-0.949) (1.092,-0.954) (1.053,-0.955) (1.014,-0.955) (0.975,-0.952) (0.936,-0.946) (0.898,-0.937) (0.860,-0.925) (0.824,-0.911) (0.788,-0.894) (0.754,-0.875) (0.722,-0.853) (0.691,-0.828) (0.662,-0.802) (0.635,-0.773) (0.622,-0.757)};
\draw[flinea, flecha=1.359cm] plot coordinates {(2.131,0.064) (2.117,0.100) (2.102,0.137) (2.088,0.173) (2.073,0.210) (2.058,0.246) (2.044,0.282) (2.028,0.318) (2.013,0.354) (1.997,0.390) (1.981,0.426) (1.965,0.462) (1.948,0.497) (1.931,0.532) (1.913,0.567) (1.895,0.602) (1.876,0.636) (1.857,0.671) (1.837,0.704) (1.817,0.738) (1.796,0.771) (1.774,0.804) (1.752,0.836) (1.729,0.868) (1.706,0.899) (1.681,0.930) (1.656,0.960) (1.630,0.990) (1.604,1.019) (1.577,1.047) (1.549,1.074) (1.520,1.101) (1.490,1.126) (1.460,1.151) (1.429,1.175) (1.397,1.198) (1.364,1.220) (1.331,1.240) (1.297,1.260) (1.262,1.278) (1.227,1.295) (1.191,1.310) (1.154,1.324) (1.117,1.336) (1.079,1.347) (1.041,1.356) (1.003,1.363) (0.964,1.369) (0.925,1.372) (0.886,1.373) (0.846,1.372) (0.807,1.369) (0.768,1.364) (0.730,1.357) (0.692,1.347) (0.655,1.334) (0.618,1.320) (0.583,1.303) (0.549,1.284) (0.516,1.262) (0.485,1.239) (0.455,1.213) (0.428,1.185) (0.402,1.156) (0.378,1.125) (0.356,1.092) (0.336,1.058) (0.319,1.023) (0.304,0.987) (0.290,0.950) (0.286,0.937)};
\draw[flinea, flecha=1.359cm] plot coordinates {(2.131,-0.064) (2.117,-0.100) (2.102,-0.137) (2.088,-0.173) (2.073,-0.210) (2.058,-0.246) (2.044,-0.282) (2.028,-0.318) (2.013,-0.354) (1.997,-0.390) (1.981,-0.426) (1.965,-0.462) (1.948,-0.497) (1.931,-0.532) (1.913,-0.567) (1.895,-0.602) (1.876,-0.636) (1.857,-0.671) (1.837,-0.704) (1.817,-0.738) (1.796,-0.771) (1.774,-0.804) (1.752,-0.836) (1.729,-0.868) (1.706,-0.899) (1.681,-0.930) (1.656,-0.960) (1.630,-0.990) (1.604,-1.019) (1.577,-1.047) (1.549,-1.074) (1.520,-1.101) (1.490,-1.126) (1.460,-1.151) (1.429,-1.175) (1.397,-1.198) (1.364,-1.220) (1.331,-1.240) (1.297,-1.260) (1.262,-1.278) (1.227,-1.295) (1.191,-1.310) (1.154,-1.324) (1.117,-1.336) (1.079,-1.347) (1.041,-1.356) (1.003,-1.363) (0.964,-1.369) (0.925,-1.372) (0.886,-1.373) (0.846,-1.372) (0.807,-1.369) (0.768,-1.364) (0.730,-1.357) (0.692,-1.347) (0.655,-1.334) (0.618,-1.320) (0.583,-1.303) (0.549,-1.284) (0.516,-1.262) (0.485,-1.239) (0.455,-1.213) (0.428,-1.185) (0.402,-1.156) (0.378,-1.125) (0.356,-1.092) (0.336,-1.058) (0.319,-1.023) (0.304,-0.987) (0.290,-0.950) (0.286,-0.937)};
\draw[flinea, flecha=2.301cm] plot coordinates {(2.144,0.067) (2.136,0.106) (2.129,0.145) (2.122,0.183) (2.114,0.221) (2.107,0.260) (2.099,0.298) (2.091,0.337) (2.082,0.375) (2.074,0.413) (2.065,0.451) (2.055,0.490) (2.046,0.528) (2.036,0.565) (2.026,0.603) (2.015,0.641) (2.004,0.678) (1.992,0.716) (1.980,0.753) (1.967,0.790) (1.954,0.827) (1.940,0.864) (1.926,0.900) (1.911,0.937) (1.896,0.973) (1.880,1.009) (1.863,1.044) (1.846,1.079) (1.828,1.114) (1.810,1.149) (1.791,1.183) (1.772,1.217) (1.751,1.251) (1.731,1.284) (1.709,1.317) (1.687,1.349) (1.664,1.381) (1.641,1.413) (1.617,1.444) (1.593,1.474) (1.568,1.505) (1.542,1.534) (1.515,1.563) (1.488,1.591) (1.461,1.619) (1.433,1.646) (1.404,1.673) (1.374,1.699) (1.344,1.724) (1.314,1.749) (1.283,1.773) (1.251,1.796) (1.219,1.818) (1.186,1.840) (1.153,1.861) (1.119,1.881) (1.085,1.900) (1.051,1.918) (1.015,1.936) (0.980,1.952) (0.944,1.968) (0.908,1.983) (0.871,1.996) (0.834,2.009) (0.796,2.021) (0.759,2.031) (0.721,2.041) (0.682,2.049) (0.644,2.057) (0.605,2.063) (0.566,2.068) (0.527,2.072) (0.488,2.074) (0.449,2.075) (0.410,2.075) (0.371,2.074) (0.331,2.071) (0.292,2.067) (0.254,2.062) (0.215,2.055) (0.177,2.047) (0.139,2.037) (0.101,2.026) (0.064,2.013) (0.028,1.998) (-0.008,1.982) (-0.043,1.965) (-0.077,1.946) (-0.111,1.925) (-0.143,1.903) (-0.174,1.880) (-0.204,1.855) (-0.233,1.828) (-0.261,1.800) (-0.287,1.770) (-0.311,1.740) (-0.333,1.708) (-0.354,1.674) (-0.373,1.640) (-0.391,1.605) (-0.406,1.569) (-0.419,1.532) (-0.431,1.495) (-0.440,1.456) (-0.447,1.418) (-0.453,1.379) (-0.456,1.340) (-0.458,1.301) (-0.458,1.262) (-0.456,1.223) (-0.452,1.184) (-0.446,1.145) (-0.439,1.106) (-0.431,1.068) (-0.421,1.030) (-0.409,0.993) (-0.397,0.955) (-0.383,0.919) (-0.377,0.905)};
\draw[flinea, flecha=2.301cm] plot coordinates {(2.144,-0.067) (2.136,-0.106) (2.129,-0.145) (2.122,-0.183) (2.114,-0.221) (2.107,-0.260) (2.099,-0.298) (2.091,-0.337) (2.082,-0.375) (2.074,-0.413) (2.065,-0.451) (2.055,-0.490) (2.046,-0.528) (2.036,-0.565) (2.026,-0.603) (2.015,-0.641) (2.004,-0.678) (1.992,-0.716) (1.980,-0.753) (1.967,-0.790) (1.954,-0.827) (1.940,-0.864) (1.926,-0.900) (1.911,-0.937) (1.896,-0.973) (1.880,-1.009) (1.863,-1.044) (1.846,-1.079) (1.828,-1.114) (1.810,-1.149) (1.791,-1.183) (1.772,-1.217) (1.751,-1.251) (1.731,-1.284) (1.709,-1.317) (1.687,-1.349) (1.664,-1.381) (1.641,-1.413) (1.617,-1.444) (1.593,-1.474) (1.568,-1.505) (1.542,-1.534) (1.515,-1.563) (1.488,-1.591) (1.461,-1.619) (1.433,-1.646) (1.404,-1.673) (1.374,-1.699) (1.344,-1.724) (1.314,-1.749) (1.283,-1.773) (1.251,-1.796) (1.219,-1.818) (1.186,-1.840) (1.153,-1.861) (1.119,-1.881) (1.085,-1.900) (1.051,-1.918) (1.015,-1.936) (0.980,-1.952) (0.944,-1.968) (0.908,-1.983) (0.871,-1.996) (0.834,-2.009) (0.796,-2.021) (0.759,-2.031) (0.721,-2.041) (0.682,-2.049) (0.644,-2.057) (0.605,-2.063) (0.566,-2.068) (0.527,-2.072) (0.488,-2.074) (0.449,-2.075) (0.410,-2.075) (0.371,-2.074) (0.331,-2.071) (0.292,-2.067) (0.254,-2.062) (0.215,-2.055) (0.177,-2.047) (0.139,-2.037) (0.101,-2.026) (0.064,-2.013) (0.028,-1.998) (-0.008,-1.982) (-0.043,-1.965) (-0.077,-1.946) (-0.111,-1.925) (-0.143,-1.903) (-0.174,-1.880) (-0.204,-1.855) (-0.233,-1.828) (-0.261,-1.800) (-0.287,-1.770) (-0.311,-1.740) (-0.333,-1.708) (-0.354,-1.674) (-0.373,-1.640) (-0.391,-1.605) (-0.406,-1.569) (-0.419,-1.532) (-0.431,-1.495) (-0.440,-1.456) (-0.447,-1.418) (-0.453,-1.379) (-0.456,-1.340) (-0.458,-1.301) (-0.458,-1.262) (-0.456,-1.223) (-0.452,-1.184) (-0.446,-1.145) (-0.439,-1.106) (-0.431,-1.068) (-0.421,-1.030) (-0.409,-0.993) (-0.397,-0.955) (-0.383,-0.919) (-0.377,-0.905)};
\draw[flinea, flecha=1.186cm] plot coordinates {(2.156,0.069) (2.156,0.108) (2.156,0.147) (2.156,0.186) (2.155,0.225) (2.155,0.265) (2.154,0.304) (2.153,0.343) (2.152,0.382) (2.151,0.421) (2.149,0.461) (2.147,0.500) (2.145,0.539) (2.142,0.578) (2.139,0.617) (2.136,0.656) (2.132,0.695) (2.128,0.734) (2.124,0.773) (2.119,0.812) (2.114,0.851) (2.108,0.890) (2.102,0.928) (2.096,0.967) (2.089,1.006) (2.082,1.044) (2.074,1.083) (2.066,1.121) (2.057,1.159) (2.048,1.197) (2.038,1.235) (2.028,1.273) (2.018,1.311) (2.006,1.348) (1.995,1.386) (1.983,1.423) (1.971,1.460) (1.958,1.497) (1.945,1.534) (1.931,1.571) (1.917,1.608) (1.902,1.644) (1.887,1.680) (1.871,1.716) (1.855,1.752) (1.839,1.787) (1.822,1.823) (1.805,1.858) (1.787,1.893) (1.769,1.928) (1.750,1.962) (1.731,1.997) (1.712,2.031) (1.692,2.065) (1.672,2.098) (1.651,2.132) (1.631,2.165) (1.609,2.198) (1.588,2.230) (1.565,2.263) (1.543,2.295) (1.532,2.311)};
\draw[flinea, flecha=1.186cm] plot coordinates {(2.156,-0.069) (2.156,-0.108) (2.156,-0.147) (2.156,-0.186) (2.155,-0.225) (2.155,-0.265) (2.154,-0.304) (2.153,-0.343) (2.152,-0.382) (2.151,-0.421) (2.149,-0.461) (2.147,-0.500) (2.145,-0.539) (2.142,-0.578) (2.139,-0.617) (2.136,-0.656) (2.132,-0.695) (2.128,-0.734) (2.124,-0.773) (2.119,-0.812) (2.114,-0.851) (2.108,-0.890) (2.102,-0.928) (2.096,-0.967) (2.089,-1.006) (2.082,-1.044) (2.074,-1.083) (2.066,-1.121) (2.057,-1.159) (2.048,-1.197) (2.038,-1.235) (2.028,-1.273) (2.018,-1.311) (2.006,-1.348) (1.995,-1.386) (1.983,-1.423) (1.971,-1.460) (1.958,-1.497) (1.945,-1.534) (1.931,-1.571) (1.917,-1.608) (1.902,-1.644) (1.887,-1.680) (1.871,-1.716) (1.855,-1.752) (1.839,-1.787) (1.822,-1.823) (1.805,-1.858) (1.787,-1.893) (1.769,-1.928) (1.750,-1.962) (1.731,-1.997) (1.712,-2.031) (1.692,-2.065) (1.672,-2.098) (1.651,-2.132) (1.631,-2.165) (1.609,-2.198) (1.588,-2.230) (1.565,-2.263) (1.543,-2.295) (1.532,-2.311)};
\draw[flinea, flecha=1.132cm] plot coordinates {(2.168,0.067) (2.176,0.106) (2.183,0.145) (2.189,0.183) (2.196,0.222) (2.203,0.260) (2.209,0.299) (2.216,0.338) (2.222,0.376) (2.228,0.415) (2.233,0.454) (2.239,0.493) (2.244,0.532) (2.249,0.571) (2.254,0.609) (2.258,0.648) (2.262,0.687) (2.266,0.726) (2.270,0.765) (2.273,0.805) (2.276,0.844) (2.278,0.883) (2.280,0.922) (2.282,0.961) (2.284,1.000) (2.285,1.039) (2.286,1.079) (2.287,1.118) (2.287,1.157) (2.287,1.196) (2.287,1.235) (2.286,1.275) (2.285,1.314) (2.283,1.353) (2.282,1.392) (2.279,1.431) (2.277,1.470) (2.274,1.509) (2.271,1.549) (2.268,1.588) (2.264,1.627) (2.260,1.666) (2.256,1.705) (2.251,1.744) (2.246,1.782) (2.241,1.821) (2.236,1.860) (2.230,1.899) (2.224,1.938) (2.217,1.976) (2.211,2.015) (2.204,2.053) (2.197,2.092) (2.189,2.130) (2.181,2.169) (2.173,2.207) (2.165,2.246) (2.156,2.284) (2.150,2.312)};
\draw[flinea, flecha=1.132cm] plot coordinates {(2.168,-0.067) (2.176,-0.106) (2.183,-0.145) (2.189,-0.183) (2.196,-0.222) (2.203,-0.260) (2.209,-0.299) (2.216,-0.338) (2.222,-0.376) (2.228,-0.415) (2.233,-0.454) (2.239,-0.493) (2.244,-0.532) (2.249,-0.571) (2.254,-0.609) (2.258,-0.648) (2.262,-0.687) (2.266,-0.726) (2.270,-0.765) (2.273,-0.805) (2.276,-0.844) (2.278,-0.883) (2.280,-0.922) (2.282,-0.961) (2.284,-1.000) (2.285,-1.039) (2.286,-1.079) (2.287,-1.118) (2.287,-1.157) (2.287,-1.196) (2.287,-1.235) (2.286,-1.275) (2.285,-1.314) (2.283,-1.353) (2.282,-1.392) (2.279,-1.431) (2.277,-1.470) (2.274,-1.509) (2.271,-1.549) (2.268,-1.588) (2.264,-1.627) (2.260,-1.666) (2.256,-1.705) (2.251,-1.744) (2.246,-1.782) (2.241,-1.821) (2.236,-1.860) (2.230,-1.899) (2.224,-1.938) (2.217,-1.976) (2.211,-2.015) (2.204,-2.053) (2.197,-2.092) (2.189,-2.130) (2.181,-2.169) (2.173,-2.207) (2.165,-2.246) (2.156,-2.284) (2.150,-2.312)};
\draw[flinea, flecha=1.156cm] plot coordinates {(2.181,0.064) (2.195,0.100) (2.209,0.137) (2.223,0.174) (2.237,0.210) (2.251,0.247) (2.265,0.284) (2.279,0.320) (2.292,0.357) (2.305,0.394) (2.318,0.431) (2.331,0.468) (2.344,0.505) (2.356,0.542) (2.369,0.579) (2.381,0.617) (2.393,0.654) (2.404,0.692) (2.416,0.729) (2.427,0.767) (2.438,0.804) (2.449,0.842) (2.459,0.880) (2.470,0.918) (2.480,0.955) (2.489,0.993) (2.499,1.031) (2.508,1.069) (2.517,1.108) (2.526,1.146) (2.535,1.184) (2.543,1.222) (2.552,1.261) (2.560,1.299) (2.567,1.337) (2.575,1.376) (2.582,1.414) (2.589,1.453) (2.596,1.492) (2.603,1.530) (2.609,1.569) (2.615,1.608) (2.621,1.646) (2.627,1.685) (2.632,1.724) (2.638,1.763) (2.643,1.802) (2.648,1.841) (2.653,1.879) (2.657,1.918) (2.662,1.957) (2.666,1.996) (2.670,2.035) (2.674,2.074) (2.677,2.113) (2.681,2.152) (2.684,2.191) (2.687,2.231) (2.690,2.270) (2.693,2.309)};
\draw[flinea, flecha=1.156cm] plot coordinates {(2.181,-0.064) (2.195,-0.100) (2.209,-0.137) (2.223,-0.174) (2.237,-0.210) (2.251,-0.247) (2.265,-0.284) (2.279,-0.320) (2.292,-0.357) (2.305,-0.394) (2.318,-0.431) (2.331,-0.468) (2.344,-0.505) (2.356,-0.542) (2.369,-0.579) (2.381,-0.617) (2.393,-0.654) (2.404,-0.692) (2.416,-0.729) (2.427,-0.767) (2.438,-0.804) (2.449,-0.842) (2.459,-0.880) (2.470,-0.918) (2.480,-0.955) (2.489,-0.993) (2.499,-1.031) (2.508,-1.069) (2.517,-1.108) (2.526,-1.146) (2.535,-1.184) (2.543,-1.222) (2.552,-1.261) (2.560,-1.299) (2.567,-1.337) (2.575,-1.376) (2.582,-1.414) (2.589,-1.453) (2.596,-1.492) (2.603,-1.530) (2.609,-1.569) (2.615,-1.608) (2.621,-1.646) (2.627,-1.685) (2.632,-1.724) (2.638,-1.763) (2.643,-1.802) (2.648,-1.841) (2.653,-1.879) (2.657,-1.918) (2.662,-1.957) (2.666,-1.996) (2.670,-2.035) (2.674,-2.074) (2.677,-2.113) (2.681,-2.152) (2.684,-2.191) (2.687,-2.231) (2.690,-2.270) (2.693,-2.309)};
\draw[flinea, flecha=1.259cm] plot coordinates {(2.193,0.057) (2.215,0.090) (2.236,0.123) (2.257,0.156) (2.278,0.189) (2.300,0.222) (2.321,0.255) (2.341,0.289) (2.362,0.322) (2.383,0.355) (2.403,0.389) (2.424,0.422) (2.444,0.456) (2.464,0.489) (2.484,0.523) (2.504,0.557) (2.524,0.591) (2.543,0.625) (2.563,0.659) (2.582,0.693) (2.601,0.727) (2.620,0.761) (2.639,0.796) (2.657,0.830) (2.676,0.865) (2.694,0.900) (2.712,0.934) (2.730,0.969) (2.748,1.004) (2.766,1.039) (2.783,1.074) (2.800,1.109) (2.818,1.144) (2.835,1.180) (2.852,1.215) (2.869,1.251) (2.885,1.286) (2.902,1.322) (2.918,1.357) (2.934,1.393) (2.950,1.429) (2.966,1.464) (2.982,1.500) (2.998,1.536) (3.013,1.572) (3.029,1.608) (3.044,1.644) (3.059,1.680) (3.074,1.717) (3.089,1.753) (3.104,1.789) (3.119,1.826) (3.133,1.862) (3.148,1.898) (3.162,1.935) (3.176,1.971) (3.191,2.008) (3.205,2.045) (3.219,2.081) (3.232,2.118) (3.246,2.155) (3.260,2.191) (3.273,2.228) (3.286,2.265) (3.300,2.302) (3.303,2.311)};
\draw[flinea, flecha=1.259cm] plot coordinates {(2.193,-0.057) (2.215,-0.090) (2.236,-0.123) (2.257,-0.156) (2.278,-0.189) (2.300,-0.222) (2.321,-0.255) (2.341,-0.289) (2.362,-0.322) (2.383,-0.355) (2.403,-0.389) (2.424,-0.422) (2.444,-0.456) (2.464,-0.489) (2.484,-0.523) (2.504,-0.557) (2.524,-0.591) (2.543,-0.625) (2.563,-0.659) (2.582,-0.693) (2.601,-0.727) (2.620,-0.761) (2.639,-0.796) (2.657,-0.830) (2.676,-0.865) (2.694,-0.900) (2.712,-0.934) (2.730,-0.969) (2.748,-1.004) (2.766,-1.039) (2.783,-1.074) (2.800,-1.109) (2.818,-1.144) (2.835,-1.180) (2.852,-1.215) (2.869,-1.251) (2.885,-1.286) (2.902,-1.322) (2.918,-1.357) (2.934,-1.393) (2.950,-1.429) (2.966,-1.464) (2.982,-1.500) (2.998,-1.536) (3.013,-1.572) (3.029,-1.608) (3.044,-1.644) (3.059,-1.680) (3.074,-1.717) (3.089,-1.753) (3.104,-1.789) (3.119,-1.826) (3.133,-1.862) (3.148,-1.898) (3.162,-1.935) (3.176,-1.971) (3.191,-2.008) (3.205,-2.045) (3.219,-2.081) (3.232,-2.118) (3.246,-2.155) (3.260,-2.191) (3.273,-2.228) (3.286,-2.265) (3.300,-2.302) (3.303,-2.311)};
\draw[flinea, flecha=1.504cm] plot coordinates {(2.206,0.047) (2.234,0.074) (2.263,0.101) (2.291,0.128) (2.320,0.155) (2.348,0.182) (2.376,0.209) (2.404,0.237) (2.433,0.264) (2.461,0.291) (2.489,0.319) (2.516,0.346) (2.544,0.374) (2.572,0.402) (2.600,0.429) (2.627,0.457) (2.655,0.485) (2.682,0.513) (2.709,0.541) (2.737,0.569) (2.764,0.598) (2.791,0.626) (2.818,0.655) (2.845,0.683) (2.872,0.712) (2.898,0.740) (2.925,0.769) (2.951,0.798) (2.978,0.827) (3.004,0.856) (3.031,0.885) (3.057,0.914) (3.083,0.943) (3.109,0.973) (3.135,1.002) (3.161,1.031) (3.187,1.061) (3.212,1.090) (3.238,1.120) (3.264,1.150) (3.289,1.180) (3.315,1.209) (3.340,1.239) (3.365,1.269) (3.390,1.299) (3.416,1.329) (3.441,1.360) (3.466,1.390) (3.491,1.420) (3.515,1.450) (3.540,1.481) (3.565,1.511) (3.590,1.541) (3.614,1.572) (3.639,1.603) (3.663,1.633) (3.688,1.664) (3.712,1.695) (3.736,1.725) (3.761,1.756) (3.785,1.787) (3.809,1.818) (3.833,1.849) (3.857,1.880) (3.881,1.911) (3.905,1.942) (3.929,1.973) (3.953,2.004) (3.976,2.035) (4.000,2.066) (4.024,2.098) (4.047,2.129) (4.071,2.160) (4.095,2.192) (4.118,2.223) (4.142,2.254) (4.165,2.286) (4.182,2.309)};
\draw[flinea, flecha=1.504cm] plot coordinates {(2.206,-0.047) (2.234,-0.074) (2.263,-0.101) (2.291,-0.128) (2.320,-0.155) (2.348,-0.182) (2.376,-0.209) (2.404,-0.237) (2.433,-0.264) (2.461,-0.291) (2.489,-0.319) (2.516,-0.346) (2.544,-0.374) (2.572,-0.402) (2.600,-0.429) (2.627,-0.457) (2.655,-0.485) (2.682,-0.513) (2.709,-0.541) (2.737,-0.569) (2.764,-0.598) (2.791,-0.626) (2.818,-0.655) (2.845,-0.683) (2.872,-0.712) (2.898,-0.740) (2.925,-0.769) (2.951,-0.798) (2.978,-0.827) (3.004,-0.856) (3.031,-0.885) (3.057,-0.914) (3.083,-0.943) (3.109,-0.973) (3.135,-1.002) (3.161,-1.031) (3.187,-1.061) (3.212,-1.090) (3.238,-1.120) (3.264,-1.150) (3.289,-1.180) (3.315,-1.209) (3.340,-1.239) (3.365,-1.269) (3.390,-1.299) (3.416,-1.329) (3.441,-1.360) (3.466,-1.390) (3.491,-1.420) (3.515,-1.450) (3.540,-1.481) (3.565,-1.511) (3.590,-1.541) (3.614,-1.572) (3.639,-1.603) (3.663,-1.633) (3.688,-1.664) (3.712,-1.695) (3.736,-1.725) (3.761,-1.756) (3.785,-1.787) (3.809,-1.818) (3.833,-1.849) (3.857,-1.880) (3.881,-1.911) (3.905,-1.942) (3.929,-1.973) (3.953,-2.004) (3.976,-2.035) (4.000,-2.066) (4.024,-2.098) (4.047,-2.129) (4.071,-2.160) (4.095,-2.192) (4.118,-2.223) (4.142,-2.254) (4.165,-2.286) (4.182,-2.309)};
\draw[flinea, flecha=1.122cm] plot coordinates {(2.218,0.029) (2.254,0.045) (2.290,0.061) (2.325,0.078) (2.361,0.094) (2.396,0.111) (2.432,0.127) (2.467,0.144) (2.503,0.160) (2.538,0.177) (2.574,0.194) (2.609,0.210) (2.645,0.227) (2.680,0.244) (2.715,0.261) (2.751,0.278) (2.786,0.295) (2.821,0.312) (2.857,0.329) (2.892,0.347) (2.927,0.364) (2.962,0.381) (2.997,0.399) (3.032,0.416) (3.067,0.434) (3.103,0.451) (3.138,0.469) (3.173,0.486) (3.208,0.504) (3.243,0.522) (3.277,0.540) (3.312,0.557) (3.347,0.575) (3.382,0.593) (3.417,0.611) (3.452,0.629) (3.486,0.647) (3.521,0.666) (3.556,0.684) (3.591,0.702) (3.625,0.720) (3.660,0.739) (3.695,0.757) (3.729,0.775) (3.764,0.794) (3.798,0.812) (3.833,0.831) (3.867,0.849) (3.902,0.868) (3.936,0.887) (3.971,0.905) (4.005,0.924) (4.040,0.943) (4.074,0.961) (4.109,0.980) (4.143,0.999) (4.177,1.018) (4.212,1.037) (4.220,1.042)};
\draw[flinea, flecha=1.122cm] plot coordinates {(2.218,-0.029) (2.254,-0.045) (2.290,-0.061) (2.325,-0.078) (2.361,-0.094) (2.396,-0.111) (2.432,-0.127) (2.467,-0.144) (2.503,-0.160) (2.538,-0.177) (2.574,-0.194) (2.609,-0.210) (2.645,-0.227) (2.680,-0.244) (2.715,-0.261) (2.751,-0.278) (2.786,-0.295) (2.821,-0.312) (2.857,-0.329) (2.892,-0.347) (2.927,-0.364) (2.962,-0.381) (2.997,-0.399) (3.032,-0.416) (3.067,-0.434) (3.103,-0.451) (3.138,-0.469) (3.173,-0.486) (3.208,-0.504) (3.243,-0.522) (3.277,-0.540) (3.312,-0.557) (3.347,-0.575) (3.382,-0.593) (3.417,-0.611) (3.452,-0.629) (3.486,-0.647) (3.521,-0.666) (3.556,-0.684) (3.591,-0.702) (3.625,-0.720) (3.660,-0.739) (3.695,-0.757) (3.729,-0.775) (3.764,-0.794) (3.798,-0.812) (3.833,-0.831) (3.867,-0.849) (3.902,-0.868) (3.936,-0.887) (3.971,-0.905) (4.005,-0.924) (4.040,-0.943) (4.074,-0.961) (4.109,-0.980) (4.143,-0.999) (4.177,-1.018) (4.212,-1.037) (4.220,-1.042)};

    \end{scope}
    \draw[gristexto,line width=0.6pt] (-\esc,0) -- (-1.75*\esc,0) -- (-1.75*\esc,-0.55);
    \tierra{(-1.75*\esc,-0.55)}
    \shade[ball color=conductor!90!white] (0,0) circle (\esc);
    % Generado automáticamente por generar_figuras.py; no editar.
\node[menos] at (0.805,0.250) {};
\node[menos] at (0.805,-0.250) {};
\node[menos] at (0.701,0.468) {};
\node[menos] at (0.701,-0.468) {};
\node[menos] at (0.535,0.651) {};
\node[menos] at (0.535,-0.651) {};
\node[menos] at (0.246,0.806) {};
\node[menos] at (0.246,-0.806) {};
\node[menos] at (-0.324,0.778) {};
\node[menos] at (-0.324,-0.778) {};

    \node[cargapos] at (2.2*\esc,0) {};
    \node[rotulo,anchor=south,fill=panel,inner sep=1pt,rounded corners=1pt] at (2.2*\esc,0.16) {$q$};
    \draw[azulborde,line width=0.5pt] \caja;
    \node[rotulogris,anchor=south west] at (-2.1*\esc,2.37*\esc) {(a) Problema real};
  \end{scope}
  % ---------------- (b)
  \begin{scope}[shift={(6.95,0)}]
    \fill[panel] \caja;
    \begin{scope}
      \clip \caja;
      \fill[panelclaro!60!fondo] (0,0) circle (\esc);
      % Generado automáticamente por generar_figuras.py; no editar.
\draw[equipot] plot coordinates {(-1.427,2.303) (-1.406,2.275) (-1.384,2.247) (-1.361,2.221) (-1.334,2.191) (-1.308,2.163) (-1.279,2.136) (-1.250,2.108) (-1.222,2.083) (-1.194,2.060) (-1.166,2.038) (-1.138,2.017) (-1.110,1.996) (-1.070,1.968) (-1.041,1.949) (-1.005,1.926) (-0.971,1.905) (-0.936,1.884) (-0.901,1.865) (-0.861,1.842) (-0.818,1.819) (-0.781,1.801) (-0.748,1.784) (-0.706,1.763) (-0.668,1.745) (-0.636,1.729) (-0.595,1.710) (-0.553,1.690) (-0.520,1.675) (-0.483,1.658) (-0.441,1.638) (-0.401,1.619) (-0.358,1.599) (-0.316,1.579) (-0.283,1.563) (-0.246,1.545) (-0.204,1.525) (-0.170,1.507) (-0.135,1.490) (-0.093,1.468) (-0.062,1.452) (-0.023,1.430) (0.014,1.410) (0.047,1.391) (0.086,1.368) (0.116,1.349) (0.154,1.326) (0.186,1.306) (0.218,1.284) (0.256,1.259) (0.297,1.229) (0.325,1.209) (0.353,1.188) (0.381,1.166) (0.409,1.144) (0.437,1.120) (0.465,1.096) (0.493,1.072) (0.520,1.047) (0.548,1.019) (0.576,0.991) (0.604,0.963) (0.629,0.935) (0.654,0.907) (0.678,0.879) (0.701,0.851) (0.722,0.823) (0.743,0.796) (0.771,0.756) (0.791,0.726) (0.813,0.691) (0.834,0.656) (0.855,0.619) (0.872,0.586) (0.893,0.544) (0.911,0.504) (0.925,0.470) (0.939,0.434) (0.953,0.393) (0.966,0.349) (0.978,0.307) (0.988,0.265) (0.996,0.223) (1.003,0.181) (1.008,0.142) (1.013,0.098) (1.015,0.056) (1.016,0.014) (1.016,-0.028) (1.015,-0.070) (1.011,-0.112) (1.007,-0.154) (1.001,-0.195) (0.994,-0.232) (0.985,-0.279) (0.974,-0.321) (0.962,-0.363) (0.949,-0.405) (0.934,-0.447) (0.917,-0.489) (0.899,-0.530) (0.883,-0.565) (0.865,-0.600) (0.842,-0.642) (0.818,-0.684) (0.799,-0.714) (0.773,-0.754) (0.743,-0.795) (0.722,-0.823) (0.701,-0.851) (0.678,-0.879) (0.654,-0.907) (0.629,-0.935) (0.604,-0.962) (0.577,-0.991) (0.549,-1.019) (0.520,-1.046) (0.493,-1.072) (0.465,-1.096) (0.437,-1.120) (0.409,-1.144) (0.381,-1.166) (0.353,-1.188) (0.325,-1.209) (0.297,-1.229) (0.259,-1.256) (0.228,-1.278) (0.197,-1.298) (0.158,-1.323) (0.116,-1.349) (0.086,-1.368) (0.047,-1.391) (0.014,-1.410) (-0.023,-1.430) (-0.062,-1.452) (-0.093,-1.468) (-0.135,-1.490) (-0.170,-1.507) (-0.204,-1.525) (-0.246,-1.545) (-0.283,-1.563) (-0.316,-1.579) (-0.358,-1.599) (-0.399,-1.619) (-0.430,-1.633) (-0.469,-1.651) (-0.511,-1.671) (-0.550,-1.689) (-0.595,-1.710) (-0.636,-1.729) (-0.668,-1.745) (-0.706,-1.763) (-0.748,-1.784) (-0.781,-1.801) (-0.818,-1.819) (-0.859,-1.842) (-0.901,-1.865) (-0.936,-1.884) (-0.971,-1.905) (-1.005,-1.926) (-1.041,-1.949) (-1.070,-1.968) (-1.110,-1.996) (-1.138,-2.017) (-1.166,-2.038) (-1.194,-2.060) (-1.222,-2.083) (-1.249,-2.108) (-1.277,-2.134) (-1.305,-2.161) (-1.333,-2.190) (-1.360,-2.219) (-1.384,-2.247) (-1.406,-2.275) (-1.427,-2.303)};
\draw[equipot] plot coordinates {(0.605,2.303) (0.594,2.261) (0.584,2.219) (0.576,2.183) (0.568,2.136) (0.562,2.095) (0.557,2.052) (0.554,2.010) (0.552,1.968) (0.551,1.926) (0.552,1.884) (0.554,1.842) (0.557,1.801) (0.561,1.759) (0.566,1.717) (0.573,1.675) (0.580,1.633) (0.589,1.591) (0.599,1.549) (0.609,1.507) (0.621,1.466) (0.632,1.429) (0.643,1.396) (0.657,1.354) (0.671,1.312) (0.687,1.270) (0.702,1.230) (0.716,1.195) (0.730,1.160) (0.743,1.125) (0.757,1.092) (0.771,1.058) (0.785,1.025) (0.799,0.991) (0.813,0.958) (0.827,0.925) (0.841,0.891) (0.855,0.857) (0.868,0.823) (0.883,0.787) (0.896,0.754) (0.911,0.714) (0.925,0.676) (0.937,0.642) (0.951,0.600) (0.964,0.558) (0.977,0.516) (0.989,0.475) (1.000,0.433) (1.008,0.396) (1.019,0.349) (1.027,0.307) (1.034,0.265) (1.040,0.223) (1.045,0.181) (1.049,0.140) (1.052,0.098) (1.054,0.056) (1.055,0.014) (1.055,-0.028) (1.054,-0.070) (1.052,-0.112) (1.048,-0.154) (1.044,-0.195) (1.038,-0.237) (1.032,-0.279) (1.024,-0.321) (1.016,-0.363) (1.006,-0.405) (0.996,-0.447) (0.985,-0.489) (0.973,-0.530) (0.960,-0.572) (0.946,-0.614) (0.932,-0.656) (0.917,-0.698) (0.901,-0.740) (0.885,-0.782) (0.869,-0.822) (0.855,-0.857) (0.841,-0.891) (0.827,-0.925) (0.813,-0.958) (0.799,-0.991) (0.785,-1.025) (0.771,-1.058) (0.757,-1.092) (0.743,-1.125) (0.730,-1.158) (0.716,-1.195) (0.702,-1.228) (0.688,-1.267) (0.674,-1.305) (0.661,-1.340) (0.647,-1.382) (0.634,-1.424) (0.621,-1.466) (0.609,-1.507) (0.599,-1.549) (0.590,-1.586) (0.580,-1.633) (0.573,-1.675) (0.566,-1.717) (0.561,-1.759) (0.557,-1.801) (0.554,-1.842) (0.552,-1.884) (0.551,-1.926) (0.552,-1.968) (0.554,-2.010) (0.557,-2.052) (0.562,-2.094) (0.568,-2.136) (0.575,-2.177) (0.584,-2.219) (0.594,-2.261) (0.604,-2.299) (0.605,-2.303)};
\draw[equipot] plot coordinates {(4.214,1.184) (4.194,1.214) (4.172,1.247) (4.147,1.284) (4.126,1.312) (4.105,1.340) (4.084,1.368) (4.061,1.396) (4.038,1.424) (4.014,1.452) (3.991,1.478) (3.964,1.507) (3.937,1.535) (3.909,1.563) (3.880,1.591) (3.852,1.618) (3.824,1.643) (3.796,1.668) (3.768,1.691) (3.740,1.714) (3.712,1.735) (3.684,1.757) (3.656,1.777) (3.623,1.801) (3.587,1.825) (3.545,1.852) (3.515,1.870) (3.475,1.893) (3.441,1.912) (3.406,1.931) (3.364,1.952) (3.331,1.968) (3.294,1.985) (3.252,2.003) (3.210,2.020) (3.169,2.036) (3.127,2.050) (3.085,2.064) (3.043,2.077) (3.001,2.088) (2.960,2.099) (2.923,2.108) (2.876,2.117) (2.834,2.125) (2.792,2.132) (2.751,2.138) (2.709,2.142) (2.667,2.146) (2.625,2.149) (2.583,2.151) (2.541,2.152) (2.500,2.152) (2.458,2.151) (2.416,2.149) (2.374,2.146) (2.332,2.142) (2.291,2.136) (2.249,2.130) (2.207,2.123) (2.165,2.115) (2.123,2.105) (2.082,2.094) (2.040,2.083) (1.998,2.069) (1.956,2.055) (1.914,2.039) (1.877,2.024) (1.845,2.010) (1.803,1.990) (1.761,1.968) (1.719,1.945) (1.687,1.926) (1.649,1.902) (1.608,1.873) (1.580,1.853) (1.552,1.831) (1.524,1.809) (1.496,1.784) (1.468,1.759) (1.440,1.732) (1.413,1.703) (1.388,1.675) (1.364,1.647) (1.342,1.619) (1.321,1.591) (1.291,1.549) (1.265,1.507) (1.245,1.474) (1.226,1.438) (1.205,1.396) (1.190,1.360) (1.176,1.326) (1.162,1.288) (1.148,1.246) (1.135,1.200) (1.124,1.158) (1.115,1.117) (1.108,1.075) (1.102,1.033) (1.097,0.991) (1.093,0.949) (1.090,0.907) (1.089,0.865) (1.088,0.823) (1.088,0.782) (1.088,0.740) (1.090,0.698) (1.092,0.656) (1.094,0.614) (1.097,0.572) (1.100,0.530) (1.103,0.489) (1.106,0.447) (1.109,0.405) (1.112,0.363) (1.115,0.321) (1.118,0.279) (1.121,0.237) (1.123,0.195) (1.125,0.154) (1.126,0.112) (1.127,0.070) (1.128,0.028) (1.128,-0.014) (1.127,-0.056) (1.126,-0.098) (1.125,-0.140) (1.123,-0.181) (1.121,-0.223) (1.119,-0.265) (1.116,-0.307) (1.113,-0.349) (1.110,-0.391) (1.107,-0.433) (1.104,-0.475) (1.101,-0.516) (1.098,-0.558) (1.095,-0.600) (1.092,-0.642) (1.090,-0.684) (1.089,-0.726) (1.088,-0.768) (1.088,-0.810) (1.088,-0.851) (1.090,-0.893) (1.092,-0.934) (1.095,-0.977) (1.100,-1.019) (1.106,-1.061) (1.113,-1.103) (1.120,-1.138) (1.128,-1.172) (1.139,-1.214) (1.151,-1.256) (1.165,-1.298) (1.181,-1.340) (1.199,-1.382) (1.217,-1.421) (1.233,-1.452) (1.257,-1.493) (1.282,-1.535) (1.311,-1.577) (1.331,-1.605) (1.353,-1.633) (1.376,-1.661) (1.399,-1.687) (1.426,-1.717) (1.454,-1.745) (1.482,-1.772) (1.510,-1.797) (1.538,-1.820) (1.566,-1.842) (1.594,-1.863) (1.623,-1.884) (1.663,-1.911) (1.705,-1.937) (1.735,-1.954) (1.775,-1.976) (1.815,-1.996) (1.859,-2.016) (1.900,-2.034) (1.942,-2.050) (1.984,-2.065) (2.026,-2.078) (2.068,-2.091) (2.109,-2.102) (2.151,-2.112) (2.193,-2.120) (2.235,-2.128) (2.277,-2.135) (2.318,-2.140) (2.360,-2.144) (2.402,-2.148) (2.444,-2.150) (2.486,-2.152) (2.528,-2.152) (2.569,-2.151) (2.611,-2.150) (2.653,-2.147) (2.695,-2.144) (2.737,-2.139) (2.778,-2.134) (2.820,-2.127) (2.854,-2.122) (2.890,-2.115) (2.932,-2.106) (2.974,-2.096) (3.015,-2.085) (3.057,-2.073) (3.099,-2.060) (3.141,-2.046) (3.183,-2.031) (3.224,-2.014) (3.266,-1.997) (3.301,-1.982) (3.336,-1.966) (3.378,-1.946) (3.415,-1.926) (3.447,-1.909) (3.489,-1.885) (3.531,-1.860) (3.573,-1.834) (3.602,-1.814) (3.643,-1.787) (3.670,-1.767) (3.698,-1.746) (3.726,-1.725) (3.753,-1.703) (3.782,-1.679) (3.810,-1.655) (3.838,-1.631) (3.866,-1.605) (3.893,-1.579) (3.921,-1.551) (3.949,-1.523) (3.976,-1.493) (4.002,-1.466) (4.026,-1.438) (4.050,-1.410) (4.073,-1.382) (4.095,-1.354) (4.116,-1.326) (4.137,-1.298) (4.158,-1.268) (4.185,-1.228) (4.212,-1.186) (4.214,-1.184)};
\draw[equipot] plot coordinates {(2.193,-1.298) (2.235,-1.302) (2.277,-1.304) (2.318,-1.305) (2.360,-1.304) (2.402,-1.302) (2.444,-1.298) (2.486,-1.293) (2.528,-1.286) (2.569,-1.278) (2.611,-1.268) (2.653,-1.256) (2.695,-1.243) (2.737,-1.228) (2.772,-1.214) (2.805,-1.200) (2.848,-1.180) (2.890,-1.158) (2.932,-1.135) (2.961,-1.117) (3.001,-1.090) (3.043,-1.061) (3.071,-1.039) (3.099,-1.017) (3.127,-0.993) (3.155,-0.967) (3.183,-0.940) (3.210,-0.912) (3.238,-0.881) (3.264,-0.851) (3.287,-0.823) (3.308,-0.796) (3.328,-0.768) (3.350,-0.736) (3.374,-0.698) (3.399,-0.656) (3.420,-0.618) (3.436,-0.586) (3.455,-0.544) (3.473,-0.502) (3.489,-0.461) (3.503,-0.421) (3.517,-0.377) (3.529,-0.335) (3.539,-0.293) (3.547,-0.251) (3.555,-0.209) (3.561,-0.167) (3.565,-0.126) (3.569,-0.084) (3.570,-0.042) (3.571,0.000) (3.570,0.042) (3.569,0.084) (3.565,0.126) (3.561,0.167) (3.555,0.209) (3.547,0.251) (3.539,0.293) (3.529,0.335) (3.517,0.376) (3.504,0.419) (3.489,0.461) (3.475,0.497) (3.462,0.530) (3.447,0.562) (3.429,0.600) (3.407,0.642) (3.383,0.684) (3.364,0.715) (3.338,0.754) (3.308,0.796) (3.287,0.823) (3.264,0.851) (3.240,0.879) (3.215,0.907) (3.188,0.935) (3.159,0.963) (3.129,0.991) (3.099,1.017) (3.071,1.039) (3.043,1.061) (3.015,1.081) (2.983,1.103) (2.946,1.126) (2.915,1.145) (2.876,1.166) (2.835,1.186) (2.792,1.206) (2.751,1.223) (2.709,1.238) (2.667,1.252) (2.625,1.264) (2.583,1.275) (2.541,1.284) (2.500,1.291) (2.458,1.297) (2.416,1.301) (2.374,1.304) (2.332,1.305) (2.291,1.305) (2.249,1.303) (2.207,1.300) (2.165,1.295) (2.123,1.288) (2.082,1.280) (2.041,1.270) (1.998,1.258) (1.956,1.244) (1.914,1.229) (1.880,1.214) (1.845,1.198) (1.803,1.176) (1.772,1.158) (1.733,1.134) (1.691,1.105) (1.663,1.084) (1.636,1.061) (1.608,1.036) (1.580,1.009) (1.552,0.980) (1.525,0.949) (1.503,0.921) (1.482,0.893) (1.454,0.853) (1.436,0.823) (1.413,0.782) (1.391,0.740) (1.372,0.698) (1.357,0.659) (1.343,0.620) (1.332,0.586) (1.320,0.544) (1.310,0.502) (1.301,0.463) (1.293,0.419) (1.287,0.384) (1.282,0.349) (1.277,0.307) (1.273,0.269) (1.269,0.223) (1.267,0.181) (1.265,0.140) (1.263,0.098) (1.262,0.056) (1.262,0.014) (1.262,-0.028) (1.262,-0.070) (1.263,-0.112) (1.265,-0.154) (1.267,-0.195) (1.270,-0.237) (1.274,-0.279) (1.279,-0.321) (1.284,-0.363) (1.290,-0.405) (1.298,-0.447) (1.306,-0.489) (1.315,-0.525) (1.324,-0.558) (1.336,-0.600) (1.351,-0.642) (1.367,-0.684) (1.385,-0.726) (1.405,-0.768) (1.426,-0.807) (1.444,-0.837) (1.468,-0.874) (1.492,-0.907) (1.514,-0.935) (1.537,-0.963) (1.562,-0.991) (1.590,-1.019) (1.620,-1.047) (1.649,-1.072) (1.677,-1.095) (1.705,-1.115) (1.747,-1.143) (1.789,-1.168) (1.822,-1.186) (1.859,-1.204) (1.900,-1.223) (1.942,-1.239) (1.984,-1.254) (2.026,-1.266) (2.068,-1.277) (2.102,-1.284) (2.137,-1.290) (2.179,-1.296) (2.193,-1.298)};
\draw[equipot] plot coordinates {(2.082,-0.783) (2.123,-0.790) (2.165,-0.793) (2.207,-0.795) (2.249,-0.794) (2.291,-0.791) (2.332,-0.785) (2.374,-0.777) (2.412,-0.768) (2.457,-0.754) (2.495,-0.740) (2.527,-0.726) (2.569,-0.705) (2.607,-0.684) (2.639,-0.663) (2.669,-0.642) (2.704,-0.614) (2.736,-0.586) (2.764,-0.558) (2.790,-0.530) (2.814,-0.502) (2.835,-0.475) (2.862,-0.435) (2.881,-0.405) (2.904,-0.363) (2.923,-0.321) (2.940,-0.279) (2.954,-0.237) (2.966,-0.195) (2.974,-0.162) (2.980,-0.126) (2.986,-0.084) (2.989,-0.042) (2.990,0.000) (2.989,0.042) (2.986,0.084) (2.980,0.126) (2.974,0.162) (2.966,0.195) (2.954,0.237) (2.940,0.279) (2.923,0.321) (2.904,0.362) (2.881,0.405) (2.862,0.435) (2.835,0.475) (2.814,0.502) (2.790,0.530) (2.764,0.558) (2.737,0.586) (2.709,0.610) (2.681,0.633) (2.653,0.654) (2.611,0.681) (2.569,0.705) (2.528,0.726) (2.495,0.740) (2.458,0.754) (2.416,0.767) (2.374,0.777) (2.332,0.785) (2.291,0.791) (2.249,0.794) (2.207,0.795) (2.165,0.793) (2.123,0.790) (2.082,0.783) (2.040,0.775) (1.998,0.763) (1.956,0.749) (1.914,0.732) (1.873,0.712) (1.831,0.688) (1.789,0.660) (1.761,0.639) (1.733,0.616) (1.705,0.590) (1.677,0.561) (1.651,0.530) (1.629,0.502) (1.608,0.471) (1.584,0.433) (1.566,0.400) (1.548,0.363) (1.530,0.321) (1.516,0.279) (1.503,0.237) (1.493,0.195) (1.485,0.154) (1.479,0.112) (1.475,0.070) (1.473,0.028) (1.473,-0.014) (1.474,-0.056) (1.477,-0.098) (1.482,-0.135) (1.490,-0.181) (1.500,-0.223) (1.510,-0.261) (1.524,-0.304) (1.538,-0.340) (1.552,-0.372) (1.568,-0.405) (1.592,-0.447) (1.619,-0.489) (1.640,-0.516) (1.663,-0.544) (1.688,-0.572) (1.716,-0.600) (1.747,-0.628) (1.775,-0.650) (1.803,-0.670) (1.845,-0.696) (1.886,-0.719) (1.928,-0.738) (1.969,-0.754) (2.012,-0.767) (2.054,-0.778) (2.082,-0.783)};
\draw[equipot] plot coordinates {(2.123,-0.476) (2.165,-0.478) (2.207,-0.477) (2.249,-0.472) (2.291,-0.463) (2.332,-0.450) (2.372,-0.433) (2.416,-0.409) (2.458,-0.381) (2.486,-0.358) (2.514,-0.330) (2.541,-0.298) (2.566,-0.265) (2.590,-0.223) (2.610,-0.181) (2.625,-0.140) (2.635,-0.098) (2.642,-0.056) (2.645,-0.014) (2.645,0.028) (2.640,0.070) (2.632,0.112) (2.620,0.154) (2.604,0.195) (2.583,0.236) (2.556,0.279) (2.534,0.307) (2.509,0.335) (2.480,0.363) (2.444,0.391) (2.402,0.417) (2.360,0.438) (2.318,0.454) (2.277,0.466) (2.235,0.474) (2.193,0.478) (2.151,0.478) (2.114,0.475) (2.068,0.467) (2.026,0.455) (1.984,0.439) (1.943,0.419) (1.900,0.392) (1.872,0.370) (1.845,0.345) (1.817,0.315) (1.789,0.279) (1.763,0.237) (1.742,0.195) (1.726,0.154) (1.714,0.112) (1.706,0.070) (1.702,0.028) (1.701,-0.014) (1.704,-0.056) (1.711,-0.098) (1.719,-0.131) (1.731,-0.167) (1.747,-0.207) (1.771,-0.251) (1.799,-0.293) (1.822,-0.321) (1.849,-0.349) (1.881,-0.377) (1.914,-0.401) (1.956,-0.426) (1.998,-0.445) (2.040,-0.459) (2.082,-0.470) (2.123,-0.476)};

      % Generado automáticamente por generar_figuras.py; no editar.
\draw[flinea, flecha=0.608cm] plot coordinates {(2.094,0.029) (2.058,0.045) (2.022,0.061) (1.987,0.077) (1.951,0.094) (1.915,0.110) (1.879,0.125) (1.843,0.141) (1.807,0.157) (1.771,0.172) (1.735,0.187) (1.699,0.202) (1.662,0.216) (1.626,0.230) (1.589,0.243) (1.552,0.256) (1.514,0.268) (1.477,0.279) (1.439,0.290) (1.401,0.299) (1.363,0.307) (1.324,0.314) (1.285,0.320) (1.246,0.324) (1.207,0.326) (1.168,0.327) (1.129,0.326) (1.090,0.323) (1.051,0.318) (1.012,0.311) (0.974,0.302) (0.936,0.291) (0.936,0.291)};
\draw[flinea, flecha=0.608cm] plot coordinates {(2.094,-0.029) (2.058,-0.045) (2.022,-0.061) (1.987,-0.077) (1.951,-0.094) (1.915,-0.110) (1.879,-0.125) (1.843,-0.141) (1.807,-0.157) (1.771,-0.172) (1.735,-0.187) (1.699,-0.202) (1.662,-0.216) (1.626,-0.230) (1.589,-0.243) (1.552,-0.256) (1.514,-0.268) (1.477,-0.279) (1.439,-0.290) (1.401,-0.299) (1.363,-0.307) (1.324,-0.314) (1.285,-0.320) (1.246,-0.324) (1.207,-0.326) (1.168,-0.327) (1.129,-0.326) (1.090,-0.323) (1.051,-0.318) (1.012,-0.311) (0.974,-0.302) (0.936,-0.291) (0.936,-0.291)};
\draw[flinea, flecha=0.751cm] plot coordinates {(2.106,0.047) (2.078,0.074) (2.049,0.101) (2.020,0.128) (1.992,0.154) (1.963,0.181) (1.934,0.207) (1.905,0.233) (1.875,0.259) (1.846,0.285) (1.816,0.311) (1.786,0.336) (1.755,0.360) (1.725,0.385) (1.693,0.408) (1.662,0.432) (1.630,0.454) (1.597,0.476) (1.564,0.497) (1.530,0.517) (1.496,0.535) (1.461,0.553) (1.425,0.570) (1.389,0.585) (1.352,0.598) (1.315,0.610) (1.277,0.620) (1.239,0.628) (1.200,0.634) (1.161,0.637) (1.122,0.638) (1.082,0.637) (1.043,0.632) (1.005,0.626) (0.967,0.616) (0.930,0.604) (0.893,0.589) (0.858,0.571) (0.825,0.551) (0.815,0.545)};
\draw[flinea, flecha=0.751cm] plot coordinates {(2.106,-0.047) (2.078,-0.074) (2.049,-0.101) (2.020,-0.128) (1.992,-0.154) (1.963,-0.181) (1.934,-0.207) (1.905,-0.233) (1.875,-0.259) (1.846,-0.285) (1.816,-0.311) (1.786,-0.336) (1.755,-0.360) (1.725,-0.385) (1.693,-0.408) (1.662,-0.432) (1.630,-0.454) (1.597,-0.476) (1.564,-0.497) (1.530,-0.517) (1.496,-0.535) (1.461,-0.553) (1.425,-0.570) (1.389,-0.585) (1.352,-0.598) (1.315,-0.610) (1.277,-0.620) (1.239,-0.628) (1.200,-0.634) (1.161,-0.637) (1.122,-0.638) (1.082,-0.637) (1.043,-0.632) (1.005,-0.626) (0.967,-0.616) (0.930,-0.604) (0.893,-0.589) (0.858,-0.571) (0.825,-0.551) (0.815,-0.545)};
\draw[flinea, flecha=0.971cm] plot coordinates {(2.119,0.057) (2.097,0.090) (2.076,0.123) (2.054,0.156) (2.032,0.188) (2.011,0.221) (1.989,0.253) (1.966,0.286) (1.944,0.318) (1.921,0.350) (1.898,0.382) (1.875,0.413) (1.851,0.444) (1.827,0.475) (1.802,0.505) (1.777,0.535) (1.751,0.565) (1.725,0.594) (1.698,0.623) (1.670,0.650) (1.642,0.678) (1.613,0.704) (1.584,0.730) (1.553,0.755) (1.522,0.779) (1.490,0.801) (1.458,0.823) (1.424,0.843) (1.390,0.863) (1.355,0.880) (1.319,0.896) (1.283,0.911) (1.246,0.923) (1.208,0.934) (1.170,0.943) (1.131,0.949) (1.092,0.954) (1.053,0.955) (1.014,0.955) (0.975,0.952) (0.936,0.946) (0.898,0.937) (0.860,0.925) (0.824,0.911) (0.788,0.894) (0.754,0.875) (0.722,0.853) (0.691,0.828) (0.662,0.802) (0.635,0.773) (0.622,0.757)};
\draw[flinea, flecha=0.971cm] plot coordinates {(2.119,-0.057) (2.097,-0.090) (2.076,-0.123) (2.054,-0.156) (2.032,-0.188) (2.011,-0.221) (1.989,-0.253) (1.966,-0.286) (1.944,-0.318) (1.921,-0.350) (1.898,-0.382) (1.875,-0.413) (1.851,-0.444) (1.827,-0.475) (1.802,-0.505) (1.777,-0.535) (1.751,-0.565) (1.725,-0.594) (1.698,-0.623) (1.670,-0.650) (1.642,-0.678) (1.613,-0.704) (1.584,-0.730) (1.553,-0.755) (1.522,-0.779) (1.490,-0.801) (1.458,-0.823) (1.424,-0.843) (1.390,-0.863) (1.355,-0.880) (1.319,-0.896) (1.283,-0.911) (1.246,-0.923) (1.208,-0.934) (1.170,-0.943) (1.131,-0.949) (1.092,-0.954) (1.053,-0.955) (1.014,-0.955) (0.975,-0.952) (0.936,-0.946) (0.898,-0.937) (0.860,-0.925) (0.824,-0.911) (0.788,-0.894) (0.754,-0.875) (0.722,-0.853) (0.691,-0.828) (0.662,-0.802) (0.635,-0.773) (0.622,-0.757)};
\draw[flinea, flecha=1.359cm] plot coordinates {(2.131,0.064) (2.117,0.100) (2.102,0.137) (2.088,0.173) (2.073,0.210) (2.058,0.246) (2.044,0.282) (2.028,0.318) (2.013,0.354) (1.997,0.390) (1.981,0.426) (1.965,0.462) (1.948,0.497) (1.931,0.532) (1.913,0.567) (1.895,0.602) (1.876,0.636) (1.857,0.671) (1.837,0.704) (1.817,0.738) (1.796,0.771) (1.774,0.804) (1.752,0.836) (1.729,0.868) (1.706,0.899) (1.681,0.930) (1.656,0.960) (1.630,0.990) (1.604,1.019) (1.577,1.047) (1.549,1.074) (1.520,1.101) (1.490,1.126) (1.460,1.151) (1.429,1.175) (1.397,1.198) (1.364,1.220) (1.331,1.240) (1.297,1.260) (1.262,1.278) (1.227,1.295) (1.191,1.310) (1.154,1.324) (1.117,1.336) (1.079,1.347) (1.041,1.356) (1.003,1.363) (0.964,1.369) (0.925,1.372) (0.886,1.373) (0.846,1.372) (0.807,1.369) (0.768,1.364) (0.730,1.357) (0.692,1.347) (0.655,1.334) (0.618,1.320) (0.583,1.303) (0.549,1.284) (0.516,1.262) (0.485,1.239) (0.455,1.213) (0.428,1.185) (0.402,1.156) (0.378,1.125) (0.356,1.092) (0.336,1.058) (0.319,1.023) (0.304,0.987) (0.290,0.950) (0.286,0.937)};
\draw[flinea, flecha=1.359cm] plot coordinates {(2.131,-0.064) (2.117,-0.100) (2.102,-0.137) (2.088,-0.173) (2.073,-0.210) (2.058,-0.246) (2.044,-0.282) (2.028,-0.318) (2.013,-0.354) (1.997,-0.390) (1.981,-0.426) (1.965,-0.462) (1.948,-0.497) (1.931,-0.532) (1.913,-0.567) (1.895,-0.602) (1.876,-0.636) (1.857,-0.671) (1.837,-0.704) (1.817,-0.738) (1.796,-0.771) (1.774,-0.804) (1.752,-0.836) (1.729,-0.868) (1.706,-0.899) (1.681,-0.930) (1.656,-0.960) (1.630,-0.990) (1.604,-1.019) (1.577,-1.047) (1.549,-1.074) (1.520,-1.101) (1.490,-1.126) (1.460,-1.151) (1.429,-1.175) (1.397,-1.198) (1.364,-1.220) (1.331,-1.240) (1.297,-1.260) (1.262,-1.278) (1.227,-1.295) (1.191,-1.310) (1.154,-1.324) (1.117,-1.336) (1.079,-1.347) (1.041,-1.356) (1.003,-1.363) (0.964,-1.369) (0.925,-1.372) (0.886,-1.373) (0.846,-1.372) (0.807,-1.369) (0.768,-1.364) (0.730,-1.357) (0.692,-1.347) (0.655,-1.334) (0.618,-1.320) (0.583,-1.303) (0.549,-1.284) (0.516,-1.262) (0.485,-1.239) (0.455,-1.213) (0.428,-1.185) (0.402,-1.156) (0.378,-1.125) (0.356,-1.092) (0.336,-1.058) (0.319,-1.023) (0.304,-0.987) (0.290,-0.950) (0.286,-0.937)};
\draw[flinea, flecha=2.301cm] plot coordinates {(2.144,0.067) (2.136,0.106) (2.129,0.145) (2.122,0.183) (2.114,0.221) (2.107,0.260) (2.099,0.298) (2.091,0.337) (2.082,0.375) (2.074,0.413) (2.065,0.451) (2.055,0.490) (2.046,0.528) (2.036,0.565) (2.026,0.603) (2.015,0.641) (2.004,0.678) (1.992,0.716) (1.980,0.753) (1.967,0.790) (1.954,0.827) (1.940,0.864) (1.926,0.900) (1.911,0.937) (1.896,0.973) (1.880,1.009) (1.863,1.044) (1.846,1.079) (1.828,1.114) (1.810,1.149) (1.791,1.183) (1.772,1.217) (1.751,1.251) (1.731,1.284) (1.709,1.317) (1.687,1.349) (1.664,1.381) (1.641,1.413) (1.617,1.444) (1.593,1.474) (1.568,1.505) (1.542,1.534) (1.515,1.563) (1.488,1.591) (1.461,1.619) (1.433,1.646) (1.404,1.673) (1.374,1.699) (1.344,1.724) (1.314,1.749) (1.283,1.773) (1.251,1.796) (1.219,1.818) (1.186,1.840) (1.153,1.861) (1.119,1.881) (1.085,1.900) (1.051,1.918) (1.015,1.936) (0.980,1.952) (0.944,1.968) (0.908,1.983) (0.871,1.996) (0.834,2.009) (0.796,2.021) (0.759,2.031) (0.721,2.041) (0.682,2.049) (0.644,2.057) (0.605,2.063) (0.566,2.068) (0.527,2.072) (0.488,2.074) (0.449,2.075) (0.410,2.075) (0.371,2.074) (0.331,2.071) (0.292,2.067) (0.254,2.062) (0.215,2.055) (0.177,2.047) (0.139,2.037) (0.101,2.026) (0.064,2.013) (0.028,1.998) (-0.008,1.982) (-0.043,1.965) (-0.077,1.946) (-0.111,1.925) (-0.143,1.903) (-0.174,1.880) (-0.204,1.855) (-0.233,1.828) (-0.261,1.800) (-0.287,1.770) (-0.311,1.740) (-0.333,1.708) (-0.354,1.674) (-0.373,1.640) (-0.391,1.605) (-0.406,1.569) (-0.419,1.532) (-0.431,1.495) (-0.440,1.456) (-0.447,1.418) (-0.453,1.379) (-0.456,1.340) (-0.458,1.301) (-0.458,1.262) (-0.456,1.223) (-0.452,1.184) (-0.446,1.145) (-0.439,1.106) (-0.431,1.068) (-0.421,1.030) (-0.409,0.993) (-0.397,0.955) (-0.383,0.919) (-0.377,0.905)};
\draw[flinea, flecha=2.301cm] plot coordinates {(2.144,-0.067) (2.136,-0.106) (2.129,-0.145) (2.122,-0.183) (2.114,-0.221) (2.107,-0.260) (2.099,-0.298) (2.091,-0.337) (2.082,-0.375) (2.074,-0.413) (2.065,-0.451) (2.055,-0.490) (2.046,-0.528) (2.036,-0.565) (2.026,-0.603) (2.015,-0.641) (2.004,-0.678) (1.992,-0.716) (1.980,-0.753) (1.967,-0.790) (1.954,-0.827) (1.940,-0.864) (1.926,-0.900) (1.911,-0.937) (1.896,-0.973) (1.880,-1.009) (1.863,-1.044) (1.846,-1.079) (1.828,-1.114) (1.810,-1.149) (1.791,-1.183) (1.772,-1.217) (1.751,-1.251) (1.731,-1.284) (1.709,-1.317) (1.687,-1.349) (1.664,-1.381) (1.641,-1.413) (1.617,-1.444) (1.593,-1.474) (1.568,-1.505) (1.542,-1.534) (1.515,-1.563) (1.488,-1.591) (1.461,-1.619) (1.433,-1.646) (1.404,-1.673) (1.374,-1.699) (1.344,-1.724) (1.314,-1.749) (1.283,-1.773) (1.251,-1.796) (1.219,-1.818) (1.186,-1.840) (1.153,-1.861) (1.119,-1.881) (1.085,-1.900) (1.051,-1.918) (1.015,-1.936) (0.980,-1.952) (0.944,-1.968) (0.908,-1.983) (0.871,-1.996) (0.834,-2.009) (0.796,-2.021) (0.759,-2.031) (0.721,-2.041) (0.682,-2.049) (0.644,-2.057) (0.605,-2.063) (0.566,-2.068) (0.527,-2.072) (0.488,-2.074) (0.449,-2.075) (0.410,-2.075) (0.371,-2.074) (0.331,-2.071) (0.292,-2.067) (0.254,-2.062) (0.215,-2.055) (0.177,-2.047) (0.139,-2.037) (0.101,-2.026) (0.064,-2.013) (0.028,-1.998) (-0.008,-1.982) (-0.043,-1.965) (-0.077,-1.946) (-0.111,-1.925) (-0.143,-1.903) (-0.174,-1.880) (-0.204,-1.855) (-0.233,-1.828) (-0.261,-1.800) (-0.287,-1.770) (-0.311,-1.740) (-0.333,-1.708) (-0.354,-1.674) (-0.373,-1.640) (-0.391,-1.605) (-0.406,-1.569) (-0.419,-1.532) (-0.431,-1.495) (-0.440,-1.456) (-0.447,-1.418) (-0.453,-1.379) (-0.456,-1.340) (-0.458,-1.301) (-0.458,-1.262) (-0.456,-1.223) (-0.452,-1.184) (-0.446,-1.145) (-0.439,-1.106) (-0.431,-1.068) (-0.421,-1.030) (-0.409,-0.993) (-0.397,-0.955) (-0.383,-0.919) (-0.377,-0.905)};
\draw[flinea, flecha=1.186cm] plot coordinates {(2.156,0.069) (2.156,0.108) (2.156,0.147) (2.156,0.186) (2.155,0.225) (2.155,0.265) (2.154,0.304) (2.153,0.343) (2.152,0.382) (2.151,0.421) (2.149,0.461) (2.147,0.500) (2.145,0.539) (2.142,0.578) (2.139,0.617) (2.136,0.656) (2.132,0.695) (2.128,0.734) (2.124,0.773) (2.119,0.812) (2.114,0.851) (2.108,0.890) (2.102,0.928) (2.096,0.967) (2.089,1.006) (2.082,1.044) (2.074,1.083) (2.066,1.121) (2.057,1.159) (2.048,1.197) (2.038,1.235) (2.028,1.273) (2.018,1.311) (2.006,1.348) (1.995,1.386) (1.983,1.423) (1.971,1.460) (1.958,1.497) (1.945,1.534) (1.931,1.571) (1.917,1.608) (1.902,1.644) (1.887,1.680) (1.871,1.716) (1.855,1.752) (1.839,1.787) (1.822,1.823) (1.805,1.858) (1.787,1.893) (1.769,1.928) (1.750,1.962) (1.731,1.997) (1.712,2.031) (1.692,2.065) (1.672,2.098) (1.651,2.132) (1.631,2.165) (1.609,2.198) (1.588,2.230) (1.565,2.263) (1.543,2.295) (1.532,2.311)};
\draw[flinea, flecha=1.186cm] plot coordinates {(2.156,-0.069) (2.156,-0.108) (2.156,-0.147) (2.156,-0.186) (2.155,-0.225) (2.155,-0.265) (2.154,-0.304) (2.153,-0.343) (2.152,-0.382) (2.151,-0.421) (2.149,-0.461) (2.147,-0.500) (2.145,-0.539) (2.142,-0.578) (2.139,-0.617) (2.136,-0.656) (2.132,-0.695) (2.128,-0.734) (2.124,-0.773) (2.119,-0.812) (2.114,-0.851) (2.108,-0.890) (2.102,-0.928) (2.096,-0.967) (2.089,-1.006) (2.082,-1.044) (2.074,-1.083) (2.066,-1.121) (2.057,-1.159) (2.048,-1.197) (2.038,-1.235) (2.028,-1.273) (2.018,-1.311) (2.006,-1.348) (1.995,-1.386) (1.983,-1.423) (1.971,-1.460) (1.958,-1.497) (1.945,-1.534) (1.931,-1.571) (1.917,-1.608) (1.902,-1.644) (1.887,-1.680) (1.871,-1.716) (1.855,-1.752) (1.839,-1.787) (1.822,-1.823) (1.805,-1.858) (1.787,-1.893) (1.769,-1.928) (1.750,-1.962) (1.731,-1.997) (1.712,-2.031) (1.692,-2.065) (1.672,-2.098) (1.651,-2.132) (1.631,-2.165) (1.609,-2.198) (1.588,-2.230) (1.565,-2.263) (1.543,-2.295) (1.532,-2.311)};
\draw[flinea, flecha=1.132cm] plot coordinates {(2.168,0.067) (2.176,0.106) (2.183,0.145) (2.189,0.183) (2.196,0.222) (2.203,0.260) (2.209,0.299) (2.216,0.338) (2.222,0.376) (2.228,0.415) (2.233,0.454) (2.239,0.493) (2.244,0.532) (2.249,0.571) (2.254,0.609) (2.258,0.648) (2.262,0.687) (2.266,0.726) (2.270,0.765) (2.273,0.805) (2.276,0.844) (2.278,0.883) (2.280,0.922) (2.282,0.961) (2.284,1.000) (2.285,1.039) (2.286,1.079) (2.287,1.118) (2.287,1.157) (2.287,1.196) (2.287,1.235) (2.286,1.275) (2.285,1.314) (2.283,1.353) (2.282,1.392) (2.279,1.431) (2.277,1.470) (2.274,1.509) (2.271,1.549) (2.268,1.588) (2.264,1.627) (2.260,1.666) (2.256,1.705) (2.251,1.744) (2.246,1.782) (2.241,1.821) (2.236,1.860) (2.230,1.899) (2.224,1.938) (2.217,1.976) (2.211,2.015) (2.204,2.053) (2.197,2.092) (2.189,2.130) (2.181,2.169) (2.173,2.207) (2.165,2.246) (2.156,2.284) (2.150,2.312)};
\draw[flinea, flecha=1.132cm] plot coordinates {(2.168,-0.067) (2.176,-0.106) (2.183,-0.145) (2.189,-0.183) (2.196,-0.222) (2.203,-0.260) (2.209,-0.299) (2.216,-0.338) (2.222,-0.376) (2.228,-0.415) (2.233,-0.454) (2.239,-0.493) (2.244,-0.532) (2.249,-0.571) (2.254,-0.609) (2.258,-0.648) (2.262,-0.687) (2.266,-0.726) (2.270,-0.765) (2.273,-0.805) (2.276,-0.844) (2.278,-0.883) (2.280,-0.922) (2.282,-0.961) (2.284,-1.000) (2.285,-1.039) (2.286,-1.079) (2.287,-1.118) (2.287,-1.157) (2.287,-1.196) (2.287,-1.235) (2.286,-1.275) (2.285,-1.314) (2.283,-1.353) (2.282,-1.392) (2.279,-1.431) (2.277,-1.470) (2.274,-1.509) (2.271,-1.549) (2.268,-1.588) (2.264,-1.627) (2.260,-1.666) (2.256,-1.705) (2.251,-1.744) (2.246,-1.782) (2.241,-1.821) (2.236,-1.860) (2.230,-1.899) (2.224,-1.938) (2.217,-1.976) (2.211,-2.015) (2.204,-2.053) (2.197,-2.092) (2.189,-2.130) (2.181,-2.169) (2.173,-2.207) (2.165,-2.246) (2.156,-2.284) (2.150,-2.312)};
\draw[flinea, flecha=1.156cm] plot coordinates {(2.181,0.064) (2.195,0.100) (2.209,0.137) (2.223,0.174) (2.237,0.210) (2.251,0.247) (2.265,0.284) (2.279,0.320) (2.292,0.357) (2.305,0.394) (2.318,0.431) (2.331,0.468) (2.344,0.505) (2.356,0.542) (2.369,0.579) (2.381,0.617) (2.393,0.654) (2.404,0.692) (2.416,0.729) (2.427,0.767) (2.438,0.804) (2.449,0.842) (2.459,0.880) (2.470,0.918) (2.480,0.955) (2.489,0.993) (2.499,1.031) (2.508,1.069) (2.517,1.108) (2.526,1.146) (2.535,1.184) (2.543,1.222) (2.552,1.261) (2.560,1.299) (2.567,1.337) (2.575,1.376) (2.582,1.414) (2.589,1.453) (2.596,1.492) (2.603,1.530) (2.609,1.569) (2.615,1.608) (2.621,1.646) (2.627,1.685) (2.632,1.724) (2.638,1.763) (2.643,1.802) (2.648,1.841) (2.653,1.879) (2.657,1.918) (2.662,1.957) (2.666,1.996) (2.670,2.035) (2.674,2.074) (2.677,2.113) (2.681,2.152) (2.684,2.191) (2.687,2.231) (2.690,2.270) (2.693,2.309)};
\draw[flinea, flecha=1.156cm] plot coordinates {(2.181,-0.064) (2.195,-0.100) (2.209,-0.137) (2.223,-0.174) (2.237,-0.210) (2.251,-0.247) (2.265,-0.284) (2.279,-0.320) (2.292,-0.357) (2.305,-0.394) (2.318,-0.431) (2.331,-0.468) (2.344,-0.505) (2.356,-0.542) (2.369,-0.579) (2.381,-0.617) (2.393,-0.654) (2.404,-0.692) (2.416,-0.729) (2.427,-0.767) (2.438,-0.804) (2.449,-0.842) (2.459,-0.880) (2.470,-0.918) (2.480,-0.955) (2.489,-0.993) (2.499,-1.031) (2.508,-1.069) (2.517,-1.108) (2.526,-1.146) (2.535,-1.184) (2.543,-1.222) (2.552,-1.261) (2.560,-1.299) (2.567,-1.337) (2.575,-1.376) (2.582,-1.414) (2.589,-1.453) (2.596,-1.492) (2.603,-1.530) (2.609,-1.569) (2.615,-1.608) (2.621,-1.646) (2.627,-1.685) (2.632,-1.724) (2.638,-1.763) (2.643,-1.802) (2.648,-1.841) (2.653,-1.879) (2.657,-1.918) (2.662,-1.957) (2.666,-1.996) (2.670,-2.035) (2.674,-2.074) (2.677,-2.113) (2.681,-2.152) (2.684,-2.191) (2.687,-2.231) (2.690,-2.270) (2.693,-2.309)};
\draw[flinea, flecha=1.259cm] plot coordinates {(2.193,0.057) (2.215,0.090) (2.236,0.123) (2.257,0.156) (2.278,0.189) (2.300,0.222) (2.321,0.255) (2.341,0.289) (2.362,0.322) (2.383,0.355) (2.403,0.389) (2.424,0.422) (2.444,0.456) (2.464,0.489) (2.484,0.523) (2.504,0.557) (2.524,0.591) (2.543,0.625) (2.563,0.659) (2.582,0.693) (2.601,0.727) (2.620,0.761) (2.639,0.796) (2.657,0.830) (2.676,0.865) (2.694,0.900) (2.712,0.934) (2.730,0.969) (2.748,1.004) (2.766,1.039) (2.783,1.074) (2.800,1.109) (2.818,1.144) (2.835,1.180) (2.852,1.215) (2.869,1.251) (2.885,1.286) (2.902,1.322) (2.918,1.357) (2.934,1.393) (2.950,1.429) (2.966,1.464) (2.982,1.500) (2.998,1.536) (3.013,1.572) (3.029,1.608) (3.044,1.644) (3.059,1.680) (3.074,1.717) (3.089,1.753) (3.104,1.789) (3.119,1.826) (3.133,1.862) (3.148,1.898) (3.162,1.935) (3.176,1.971) (3.191,2.008) (3.205,2.045) (3.219,2.081) (3.232,2.118) (3.246,2.155) (3.260,2.191) (3.273,2.228) (3.286,2.265) (3.300,2.302) (3.303,2.311)};
\draw[flinea, flecha=1.259cm] plot coordinates {(2.193,-0.057) (2.215,-0.090) (2.236,-0.123) (2.257,-0.156) (2.278,-0.189) (2.300,-0.222) (2.321,-0.255) (2.341,-0.289) (2.362,-0.322) (2.383,-0.355) (2.403,-0.389) (2.424,-0.422) (2.444,-0.456) (2.464,-0.489) (2.484,-0.523) (2.504,-0.557) (2.524,-0.591) (2.543,-0.625) (2.563,-0.659) (2.582,-0.693) (2.601,-0.727) (2.620,-0.761) (2.639,-0.796) (2.657,-0.830) (2.676,-0.865) (2.694,-0.900) (2.712,-0.934) (2.730,-0.969) (2.748,-1.004) (2.766,-1.039) (2.783,-1.074) (2.800,-1.109) (2.818,-1.144) (2.835,-1.180) (2.852,-1.215) (2.869,-1.251) (2.885,-1.286) (2.902,-1.322) (2.918,-1.357) (2.934,-1.393) (2.950,-1.429) (2.966,-1.464) (2.982,-1.500) (2.998,-1.536) (3.013,-1.572) (3.029,-1.608) (3.044,-1.644) (3.059,-1.680) (3.074,-1.717) (3.089,-1.753) (3.104,-1.789) (3.119,-1.826) (3.133,-1.862) (3.148,-1.898) (3.162,-1.935) (3.176,-1.971) (3.191,-2.008) (3.205,-2.045) (3.219,-2.081) (3.232,-2.118) (3.246,-2.155) (3.260,-2.191) (3.273,-2.228) (3.286,-2.265) (3.300,-2.302) (3.303,-2.311)};
\draw[flinea, flecha=1.504cm] plot coordinates {(2.206,0.047) (2.234,0.074) (2.263,0.101) (2.291,0.128) (2.320,0.155) (2.348,0.182) (2.376,0.209) (2.404,0.237) (2.433,0.264) (2.461,0.291) (2.489,0.319) (2.516,0.346) (2.544,0.374) (2.572,0.402) (2.600,0.429) (2.627,0.457) (2.655,0.485) (2.682,0.513) (2.709,0.541) (2.737,0.569) (2.764,0.598) (2.791,0.626) (2.818,0.655) (2.845,0.683) (2.872,0.712) (2.898,0.740) (2.925,0.769) (2.951,0.798) (2.978,0.827) (3.004,0.856) (3.031,0.885) (3.057,0.914) (3.083,0.943) (3.109,0.973) (3.135,1.002) (3.161,1.031) (3.187,1.061) (3.212,1.090) (3.238,1.120) (3.264,1.150) (3.289,1.180) (3.315,1.209) (3.340,1.239) (3.365,1.269) (3.390,1.299) (3.416,1.329) (3.441,1.360) (3.466,1.390) (3.491,1.420) (3.515,1.450) (3.540,1.481) (3.565,1.511) (3.590,1.541) (3.614,1.572) (3.639,1.603) (3.663,1.633) (3.688,1.664) (3.712,1.695) (3.736,1.725) (3.761,1.756) (3.785,1.787) (3.809,1.818) (3.833,1.849) (3.857,1.880) (3.881,1.911) (3.905,1.942) (3.929,1.973) (3.953,2.004) (3.976,2.035) (4.000,2.066) (4.024,2.098) (4.047,2.129) (4.071,2.160) (4.095,2.192) (4.118,2.223) (4.142,2.254) (4.165,2.286) (4.182,2.309)};
\draw[flinea, flecha=1.504cm] plot coordinates {(2.206,-0.047) (2.234,-0.074) (2.263,-0.101) (2.291,-0.128) (2.320,-0.155) (2.348,-0.182) (2.376,-0.209) (2.404,-0.237) (2.433,-0.264) (2.461,-0.291) (2.489,-0.319) (2.516,-0.346) (2.544,-0.374) (2.572,-0.402) (2.600,-0.429) (2.627,-0.457) (2.655,-0.485) (2.682,-0.513) (2.709,-0.541) (2.737,-0.569) (2.764,-0.598) (2.791,-0.626) (2.818,-0.655) (2.845,-0.683) (2.872,-0.712) (2.898,-0.740) (2.925,-0.769) (2.951,-0.798) (2.978,-0.827) (3.004,-0.856) (3.031,-0.885) (3.057,-0.914) (3.083,-0.943) (3.109,-0.973) (3.135,-1.002) (3.161,-1.031) (3.187,-1.061) (3.212,-1.090) (3.238,-1.120) (3.264,-1.150) (3.289,-1.180) (3.315,-1.209) (3.340,-1.239) (3.365,-1.269) (3.390,-1.299) (3.416,-1.329) (3.441,-1.360) (3.466,-1.390) (3.491,-1.420) (3.515,-1.450) (3.540,-1.481) (3.565,-1.511) (3.590,-1.541) (3.614,-1.572) (3.639,-1.603) (3.663,-1.633) (3.688,-1.664) (3.712,-1.695) (3.736,-1.725) (3.761,-1.756) (3.785,-1.787) (3.809,-1.818) (3.833,-1.849) (3.857,-1.880) (3.881,-1.911) (3.905,-1.942) (3.929,-1.973) (3.953,-2.004) (3.976,-2.035) (4.000,-2.066) (4.024,-2.098) (4.047,-2.129) (4.071,-2.160) (4.095,-2.192) (4.118,-2.223) (4.142,-2.254) (4.165,-2.286) (4.182,-2.309)};
\draw[flinea, flecha=1.122cm] plot coordinates {(2.218,0.029) (2.254,0.045) (2.290,0.061) (2.325,0.078) (2.361,0.094) (2.396,0.111) (2.432,0.127) (2.467,0.144) (2.503,0.160) (2.538,0.177) (2.574,0.194) (2.609,0.210) (2.645,0.227) (2.680,0.244) (2.715,0.261) (2.751,0.278) (2.786,0.295) (2.821,0.312) (2.857,0.329) (2.892,0.347) (2.927,0.364) (2.962,0.381) (2.997,0.399) (3.032,0.416) (3.067,0.434) (3.103,0.451) (3.138,0.469) (3.173,0.486) (3.208,0.504) (3.243,0.522) (3.277,0.540) (3.312,0.557) (3.347,0.575) (3.382,0.593) (3.417,0.611) (3.452,0.629) (3.486,0.647) (3.521,0.666) (3.556,0.684) (3.591,0.702) (3.625,0.720) (3.660,0.739) (3.695,0.757) (3.729,0.775) (3.764,0.794) (3.798,0.812) (3.833,0.831) (3.867,0.849) (3.902,0.868) (3.936,0.887) (3.971,0.905) (4.005,0.924) (4.040,0.943) (4.074,0.961) (4.109,0.980) (4.143,0.999) (4.177,1.018) (4.212,1.037) (4.220,1.042)};
\draw[flinea, flecha=1.122cm] plot coordinates {(2.218,-0.029) (2.254,-0.045) (2.290,-0.061) (2.325,-0.078) (2.361,-0.094) (2.396,-0.111) (2.432,-0.127) (2.467,-0.144) (2.503,-0.160) (2.538,-0.177) (2.574,-0.194) (2.609,-0.210) (2.645,-0.227) (2.680,-0.244) (2.715,-0.261) (2.751,-0.278) (2.786,-0.295) (2.821,-0.312) (2.857,-0.329) (2.892,-0.347) (2.927,-0.364) (2.962,-0.381) (2.997,-0.399) (3.032,-0.416) (3.067,-0.434) (3.103,-0.451) (3.138,-0.469) (3.173,-0.486) (3.208,-0.504) (3.243,-0.522) (3.277,-0.540) (3.312,-0.557) (3.347,-0.575) (3.382,-0.593) (3.417,-0.611) (3.452,-0.629) (3.486,-0.647) (3.521,-0.666) (3.556,-0.684) (3.591,-0.702) (3.625,-0.720) (3.660,-0.739) (3.695,-0.757) (3.729,-0.775) (3.764,-0.794) (3.798,-0.812) (3.833,-0.831) (3.867,-0.849) (3.902,-0.868) (3.936,-0.887) (3.971,-0.905) (4.005,-0.924) (4.040,-0.943) (4.074,-0.961) (4.109,-0.980) (4.143,-0.999) (4.177,-1.018) (4.212,-1.037) (4.220,-1.042)};

      % Generado automáticamente por generar_figuras.py; no editar.
\draw[flinea, virtual] plot coordinates {(0.935,0.291) (0.898,0.278) (0.861,0.264) (0.825,0.248) (0.790,0.231) (0.755,0.213) (0.722,0.193) (0.688,0.173) (0.655,0.151) (0.623,0.129) (0.591,0.107) (0.559,0.084) (0.527,0.061) (0.496,0.038) (0.484,0.029)};
\draw[flinea, virtual] plot coordinates {(0.935,-0.291) (0.898,-0.278) (0.861,-0.264) (0.825,-0.248) (0.790,-0.231) (0.755,-0.213) (0.722,-0.193) (0.688,-0.173) (0.655,-0.151) (0.623,-0.129) (0.591,-0.107) (0.559,-0.084) (0.527,-0.061) (0.496,-0.038) (0.484,-0.029)};
\draw[flinea, virtual] plot coordinates {(0.814,0.544) (0.782,0.521) (0.752,0.496) (0.723,0.470) (0.696,0.441) (0.670,0.412) (0.646,0.381) (0.623,0.349) (0.602,0.316) (0.582,0.283) (0.562,0.248) (0.544,0.214) (0.526,0.179) (0.510,0.143) (0.493,0.108) (0.477,0.072) (0.463,0.041)};
\draw[flinea, virtual] plot coordinates {(0.814,-0.544) (0.782,-0.521) (0.752,-0.496) (0.723,-0.470) (0.696,-0.441) (0.670,-0.412) (0.646,-0.381) (0.623,-0.349) (0.602,-0.316) (0.582,-0.283) (0.562,-0.248) (0.544,-0.214) (0.526,-0.179) (0.510,-0.143) (0.493,-0.108) (0.477,-0.072) (0.463,-0.041)};
\draw[flinea, virtual] plot coordinates {(0.621,0.757) (0.597,0.726) (0.576,0.693) (0.556,0.659) (0.538,0.624) (0.522,0.588) (0.508,0.552) (0.496,0.514) (0.486,0.477) (0.477,0.439) (0.469,0.400) (0.463,0.361) (0.458,0.323) (0.454,0.284) (0.451,0.244) (0.449,0.205) (0.447,0.166) (0.446,0.127) (0.446,0.088) (0.445,0.049)};
\draw[flinea, virtual] plot coordinates {(0.621,-0.757) (0.597,-0.726) (0.576,-0.693) (0.556,-0.659) (0.538,-0.624) (0.522,-0.588) (0.508,-0.552) (0.496,-0.514) (0.486,-0.477) (0.477,-0.439) (0.469,-0.400) (0.463,-0.361) (0.458,-0.323) (0.454,-0.284) (0.451,-0.244) (0.449,-0.205) (0.447,-0.166) (0.446,-0.127) (0.446,-0.088) (0.445,-0.049)};
\draw[flinea, virtual] plot coordinates {(0.286,0.936) (0.276,0.899) (0.267,0.860) (0.261,0.822) (0.256,0.783) (0.254,0.744) (0.253,0.704) (0.253,0.665) (0.256,0.626) (0.259,0.587) (0.264,0.548) (0.271,0.509) (0.278,0.471) (0.287,0.433) (0.297,0.395) (0.307,0.357) (0.319,0.320) (0.331,0.282) (0.344,0.245) (0.358,0.209) (0.372,0.172) (0.387,0.136) (0.402,0.100) (0.418,0.064) (0.427,0.041)};
\draw[flinea, virtual] plot coordinates {(0.286,-0.936) (0.276,-0.899) (0.267,-0.860) (0.261,-0.822) (0.256,-0.783) (0.254,-0.744) (0.253,-0.704) (0.253,-0.665) (0.256,-0.626) (0.259,-0.587) (0.264,-0.548) (0.271,-0.509) (0.278,-0.471) (0.287,-0.433) (0.297,-0.395) (0.307,-0.357) (0.319,-0.320) (0.331,-0.282) (0.344,-0.245) (0.358,-0.209) (0.372,-0.172) (0.387,-0.136) (0.402,-0.100) (0.418,-0.064) (0.427,-0.041)};
\draw[flinea, virtual] plot coordinates {(-0.376,0.904) (-0.361,0.868) (-0.344,0.832) (-0.326,0.798) (-0.307,0.763) (-0.288,0.729) (-0.267,0.696) (-0.245,0.663) (-0.223,0.631) (-0.200,0.599) (-0.176,0.568) (-0.152,0.537) (-0.127,0.507) (-0.101,0.478) (-0.075,0.448) (-0.048,0.420) (-0.021,0.391) (0.007,0.364) (0.035,0.336) (0.063,0.310) (0.092,0.283) (0.121,0.257) (0.151,0.231) (0.181,0.206) (0.211,0.181) (0.242,0.156) (0.272,0.132) (0.303,0.108) (0.334,0.084) (0.365,0.060) (0.397,0.037) (0.408,0.028)};
\draw[flinea, virtual] plot coordinates {(-0.376,-0.904) (-0.361,-0.868) (-0.344,-0.832) (-0.326,-0.798) (-0.307,-0.763) (-0.288,-0.729) (-0.267,-0.696) (-0.245,-0.663) (-0.223,-0.631) (-0.200,-0.599) (-0.176,-0.568) (-0.152,-0.537) (-0.127,-0.507) (-0.101,-0.478) (-0.075,-0.448) (-0.048,-0.420) (-0.021,-0.391) (0.007,-0.364) (0.035,-0.336) (0.063,-0.310) (0.092,-0.283) (0.121,-0.257) (0.151,-0.231) (0.181,-0.206) (0.211,-0.181) (0.242,-0.156) (0.272,-0.132) (0.303,-0.108) (0.334,-0.084) (0.365,-0.060) (0.397,-0.037) (0.408,-0.028)};

    \end{scope}
    \draw[azul,line width=0.6pt,dash pattern=on 4pt off 2pt] (0,0) circle (\esc);
    \node[cargapos] at (2.2*\esc,0) {};
    \node[cargaimg] at ({\esc/2.2},0) {};
    \node[rotulo,anchor=south,fill=panel,inner sep=1pt,rounded corners=1pt] at (2.2*\esc,0.16) {$q$};
    \node[rotulo,text=azul,anchor=north,fill=panelclaro!60!fondo,inner sep=1pt,rounded corners=1pt] at ({\esc/2.2},-0.15) {$q'$};
    \node[rotulo,text=azul,fill=panel,inner sep=1.5pt] at (-1.35*\esc,-1.5*\esc) {$\varphi=0$};
    \draw[azulborde,line width=0.5pt] \caja;
    \node[rotulogris,anchor=south west] at (-2.1*\esc,2.37*\esc) {(b) Equivalente: $q$ y su imagen $q'$};
  \end{scope}
\end{tikzpicture}
